%! Convolutional Neural Nets — Unit 8, Lesson 8.2 — Group Activity
\documentclass[10pt]{article}
\usepackage{linalg-article}
\usepackage{linalg-boxes}
\newcommand{\TallMath}[1]{$\displaystyle #1\rule[-1.4em]{0pt}{3.2em}$}
\newcommand{\vv}{\mathbf{v}}
\newcommand{\bb}{\mathbf{b}}
\newcommand{\ww}{\mathbf{w}}
\newcommand{\zero}{\mathbf{0}}
\newcommand{\relu}{\mathrm{ReLU}}

\begin{document}

\pageheader{Unit 8, Lesson 8.2}{Group Activity}
\namepartnerperiod

\vspace{0.10in}

\begin{headlinebox}{blush}
\small\textbf{Scanning for a Pattern.} A \textbf{convolution} slides a small \textbf{filter} across a signal,
taking a dot product at each spot. Written as a matrix, the filter's weights \emph{repeat down the diagonals}
--- the same few weights, \textbf{reused everywhere}. That is what lets a network find the same pattern no
matter where it sits, with almost no weights. Work with your partner, show each step, and \emph{explain} what
you find.
\end{headlinebox}

\vspace{0.10in}

\begin{tcolorbox}[colback=white, colframe=black!40, title=\textbf{Tier R --- Remediate},
    fonttitle=\bfseries, arc=1.5mm, left=3mm, right=3mm, top=2mm, bottom=2mm, breakable]
Slide two different filters and see what each one does.
\begin{enumerate}[leftmargin=1.7em, itemsep=8pt]
  \item \textbf{Edge filter.} Slide $\ww=[-1,1]$ (output $=-v_i+v_{i+1}$) across $\vv=[1,1,4,4]$:
        \[ y_1=\blank{0.9cm},\ \ y_2=\blank{0.9cm},\ \ y_3=\blank{0.9cm}. \]
        Output $=[\ \blank{0.5cm},\ \blank{0.5cm},\ \blank{0.5cm}\ ]$. At which spot does it \emph{fire},
        and what does that spot correspond to in the signal? \writeline
  \item \textbf{Averaging (smoothing) filter.} Slide $\ww=\tfrac12[1,1]$ (output $=\tfrac{v_i+v_{i+1}}{2}$)
        across the same $\vv=[1,1,4,4]$:
        \[ y_1=\blank{0.9cm},\ \ y_2=\blank{0.9cm},\ \ y_3=\blank{0.9cm}. \]
        Compared to the edge filter, what does this filter \emph{do} to the jump? \writeline
\end{enumerate}
\end{tcolorbox}

\begin{tcolorbox}[colback=white, colframe=black!40, title=\textbf{Tier A --- Approaching Proficiency},
    fonttitle=\bfseries, arc=1.5mm, left=3mm, right=3mm, top=2mm, bottom=2mm, breakable]
Write the convolution as a matrix, count the weights it saves, then finish the layer.
\begin{enumerate}[leftmargin=1.7em, itemsep=9pt]
  \item \textbf{Build the matrix.} For the filter $\ww=[-1,1]$ and a length-$5$ signal, fill in the
        convolution matrix (the same two numbers slide over one step each row):
        \[ A=\begin{bmatrix} -1 & 1 & 0 & 0 & 0\\ 0 & \blank{0.5cm} & \blank{0.5cm} & 0 & 0\\
             0 & 0 & \blank{0.5cm} & \blank{0.5cm} & 0\\ 0 & 0 & 0 & \blank{0.5cm} & \blank{0.5cm} \end{bmatrix}. \]
        Then compute $A\vv$ for $\vv=[3,3,3,8,8]$: \quad output $=[\ \blank{0.5cm},\ \blank{0.5cm},\ \blank{0.5cm},\ \blank{0.5cm}\ ]$.
  \item \textbf{Count the weights.} A \emph{full} $4\times5$ layer would have how many separate weights? How
        many does this shared-weight convolution use? \writeline
  \item \textbf{Finish the layer.} Apply $\relu(A\vv+\bb)$ with bias $\bb=-2$ to the output of problem~1.
        What comes out, and which spot survives? \writeline
\end{enumerate}
\end{tcolorbox}

\begin{tcolorbox}[colback=white, colframe=black!40, title=\textbf{Tier E --- Extension},
    fonttitle=\bfseries, arc=1.5mm, left=3mm, right=3mm, top=2mm, bottom=2mm, breakable]
\textbf{The same filter, wherever the pattern hides.} Weight sharing means one filter detects its pattern
\emph{everywhere}.
\begin{enumerate}[leftmargin=1.7em, itemsep=9pt]
  \item Slide $\ww=[-1,1]$ across each signal (both have the \emph{same} jump of $6$, in different places):
        \[ \text{A: }[3,9,9,9,9]\ \to\ [\ \blank{0.5cm},\ \blank{0.5cm},\ \blank{0.5cm},\ \blank{0.5cm}\ ],
           \qquad \text{B: }[3,3,3,3,9]\ \to\ [\ \blank{0.5cm},\ \blank{0.5cm},\ \blank{0.5cm},\ \blank{0.5cm}\ ]. \]
        The same filter fired with the same strength in both --- just at a different \emph{position}. In one
        sentence, why does reusing one filter make the detector work no matter where the edge sits? \writeline
  \item A real $1000\times1000$ image has a million pixels. A full layer to a million outputs needs about
        $10^{12}$ weights; a $3\times3$ filter needs $9$. In a sentence, why is this the difference between a
        network you \emph{can} train and one you cannot? \writeline
\end{enumerate}
\end{tcolorbox}

\end{document}
